% =============================================================================
% Paper 2 R1 SEVERE 1 — Transaction Cost Sensitivity (proposed body addition)
% Suggested placement: paper/taiwan-vt/body.tex, end of §4 (immediately after
% Table~\ref{tab:vt_common} note on common-period comparison, ~line 295), OR
% as new §4.6 "Transaction Cost Sensitivity" before §5 high-frequency section.
%
% Word count: 198 (target ~150-200 EN words; deliberately compact, PBFJ-tier).
% Data source: experiments/paper2_R1_transaction_tax_fix/results.json
% Length-justified: gemini SEVERE 1 demanded explicit sensitivity, so a single
% paragraph is the minimum credible response; an inline table is included to
% give reviewers byte-traceable evidence per paper-workflow.md § Table row binding.
% =============================================================================

\subsection{Transaction Cost Sensitivity}\label{sec:tx_sensitivity}

Taiwan's Securities Transaction Tax differentiates between common stocks (0.30\% on
sell) and ETFs (0.10\% on sell), and online brokerages typically charge 30\% of the
statutory 0.1425\% commission rate. Our baseline applies 0.186\% round-trip---the
appropriate rate for 0050.TW as an ETF---but a referee may legitimately ask whether
the volatility-targeting (VT) advantage survives across the full Taiwan tax-and-commission
range. Table~\ref{tab:tx_sensitivity} reports Sharpe ratios and maximum-drawdown (MDD)
improvements over Buy \& Hold at four round-trip transaction rates spanning 0.10\% to
0.30\%. The Sharpe ratio of each strategy declines monotonically with the transaction
rate, as expected, but two findings emerge clearly. First, the MDD improvement
relative to Buy \& Hold survives the entire range: at the worst case of 0.30\%, all
four VT specifications retain MDD improvements of at least 11.5 percentage points.
Second, the monthly-rebalanced $8.63/\text{VIX}$ specification is far less sensitive to
transaction costs than the daily-rebalanced VT variants---its Sharpe range across the
0.10\%--0.30\% grid is only 0.024 versus 0.098--0.127 for the daily-rebalanced
strategies---reinforcing our practical recommendation in Section~\ref{sec:implementation}
that retail investors should adopt monthly rebalancing.

\begin{table}[!htbp]
\centering
\caption{Transaction Cost Sensitivity for VT Strategies on 0050.TW}
\label{tab:tx_sensitivity}
\small
\begin{tabular}{lccccc}
\toprule
& \multicolumn{4}{c}{Net Sharpe at round-trip transaction rate} & MDD impr.\ vs BH \\
\cmidrule(lr){2-5}
Strategy & 0.10\% & 0.15\% & 0.186\% & 0.30\% & at TX=0.30\% (pp) \\
\midrule
EWMA VT (10\%)        & 0.564 & 0.540 & 0.522 & 0.467 & $+13.6$ \\
GARCH VT (10\%)       & 0.848 & 0.817 & 0.795 & 0.725 & $+12.3$ \\
GJR VT (10\%)         & 0.955 & 0.923 & 0.900 & 0.828 & $+11.5$ \\
$8.63/\text{VIX}$ (monthly) & 0.967 & 0.961 & 0.956 & 0.943 & $+21.2$ \\
\bottomrule
\end{tabular}
% source: experiments/paper2_R1_transaction_tax_fix/results.json
%   tx_sensitivity_sweep.<strat>.{etf_floor,brief_mid,paper_canonical,stock_high}.sharpe
%   vt_vs_bh_survival.<strat>.stock_high.mdd_improvement_pp
\smallskip
\\
\small \textit{Notes:} 0.10\% is the ETF-only floor (tax 0.10\%, no commission); 0.15\%
is the midpoint of the Taiwan transaction-tax range; 0.186\% is the paper canonical (ETF
tax 0.10\% + online-discount commission 0.04275\% $\times$ 2); 0.30\% is the
hypothetical common-stock rate (not actually applicable to the 0050.TW ETF; included
for upper-bound robustness only). Evaluation windows: EWMA 2010--2026
($n=3{,}980$); GARCH/GJR 2020--2026 ($n=1{,}524$); $8.63/\text{VIX}$ 2016--2026
($n=2{,}497$). Sample annual turnovers, identical across TX rates by construction:
EWMA 489\%/yr, GARCH 634\%/yr, GJR 662\%/yr, $8.63/\text{VIX}$ 104\%/yr. MDD improvement
is in percentage points and is computed within each strategy's own evaluation window.
Reproduction: \texttt{experiments/paper2\_R1\_transaction\_tax\_fix/turnover\_and\_tax.py}
using the pinned snapshot \texttt{paper/taiwan-vt/data/0050\_tw\_twii\_2330\_tw\_2317\_tw\_2454\_tw\_0056\_tw\_spy\_vix\_2008-2026.csv}.
\end{table}
